\section{Access the labs can't withhold}
\label{sec:28.3}
The first piece is a right, written in statute. A research team the review board approves may obtain, for a stated protocol, the weights, checkpoints, update streams and release histories of any model above the compute line, and the decision records of any deployment a floor covers. What it obtains stays under the preservation custodian's security regime: the team works on the weights where the custodian holds them, and nothing leaves but results. It publishes what it finds. The operator may ask the review board to hold back a detail that would help someone steal or attack a model. The request names the detail and the attack it would enable; the burden is the operator's, as it is for any claim by the party that controls the evidence; the board answers on a fixed timetable; a request for everything is refused as one, and a request the board hasn't answered when the clock runs out is refused; and each request and its answer go on a published record. Otherwise the operator asks for everything and waits. Classified systems are reached by cleared teams working inside the enclave, as a cleared bearer works, and their findings are published by class, as notice is. What this doesn't reach is a lab outside the jurisdiction, except through the procurement term of Table~F.1a, under which a deployment the floor covers licenses only models that submit to the right.

The European Union has built part of this for platforms. The Digital Services Act lets researchers vetted by national coordinators obtain data from the largest platforms and search engines, to study the systemic risks those services pose \autocite{eu2022dsa}, and a delegated regulation of 2025 sets the procedure, with independence from commercial interests among the conditions of vetting \autocite{eu2025dsadataaccess}. It reaches data about a platform's operation, not models. What it shows is that the obligation can be written: a company that would rather not be studied can be made to answer a request it didn't invite.

Access needs something to have access to. A plurality of lineages is what tests 5 and 7 compare and what the commons is made of, and it erodes in the ordinary way markets concentrate. In January 2025 the Federal Trade Commission's staff reported on three partnerships between cloud providers and model developers, Microsoft with OpenAI, Amazon with Anthropic and Google with Anthropic. Together they involved more than \$20 billion in investment and gave the cloud providers equity and revenue-sharing rights and certain rights of consultation, control and exclusivity, and the staff flagged the risk that such arrangements limit rivals' access to computing and to engineers \autocite{ftc2025partnerships}. A lab whose compute comes from its investor has a second custodian it can't easily refuse. Partnerships like these should get the review a merger gets, and the test should be the commons' test for independence: whether the lineage's history can go a way its supplier doesn't want.

Interoperability is the rest. The Digital Markets Act requires the largest gatekeepers to let others use, free of charge, the hardware and software features their own services use \autocite{eu2022dma}. A statute here would carry that down to accelerators: firmware and drivers that let the owner of a device run a model its vendor didn't choose. Without it the home scale exists on paper, and so does the hardware a bearer is owed.
